% Pursuing Stacks (cotes 134-2 to 134-8): archive cote 134-7, site page 31 (PDF page 32), typescript page 473.
% typescript. Draft compared: ps6.tex:1438-1502.
% AI-transcribed by Claude Opus 5.5 (claude-opus-5-5, adaptive thinking, effort high), Message Batches API, batch ps3 (msgbatch_01XdFmuVVBAcgc7EFxdtCqsj), 2026-11-12; not yet checked by eye.
\scanpage{134-7/31}
\mspage{473}

For a given ground ring k,  the basic category\inserted{we'll use} of ``schematic'' objects
over k is the category of \uline{unipotent} \uline{bundles} over k,  which may be defined
as the category of \xout{xxx} functors  from Alg$_{/k}$ to \xout{sets of the type W(} (Sets)
isomorphic to functors of the type
\[ W(M) = ( k' \hand{\longmapsto} M\otimes_k k' ) \qquad , \]
where M is any k-module.  We do not restrict here M to be projective or
flat,  as we definitely want to have,  for a ring homomorphism $k \hand{\to} k'$ ,
a problemless functor  ``\uline{restriction of scalars}''
\[ U(k') \hand{\to} U(k) \qquad , \]
inserting in the commutative diagram
\diagram{
(I) \qquad
\begin{tikzcd}
Ab_{k'} \arrow[r, "\text{restr.}"] \arrow[d, "W_{k'}"'] & Ab_k \arrow[d, "W_k"] \\
U(k') \arrow[r, "\text{restr.}"] & U(k)
\end{tikzcd}
\qquad .
}
Another reason is that we want that the k-modules of the type $\hand{\pi}_i(X_{\hand{*}})$
which will come out should be eligible for defining objects in U(k).
We are more specifically interested though in ``\uline{smooth}'' objects of U(k),
namely those which correspond to \uline{projective} k-modules.  \hand{(}We prefer to
call them ``smooth'' rather than ``projective'', in order to avoid confusion
with the notion of a projective object in the usual categorical sense
for U(k).)\hand{\unsure{)}}  Another relevant notion is the notion of a \uline{submersion} ,
namely a map in U(k) \inserted{isomorphic to one} coming from an epimorphism \xout{in} $M \hand{\to} N$  in Ab$_k$ .

(If the latter can be chosen to have a projective kernel,  we may speak
of a smooth submersion.)  The ring restriction functor transforms sub-%
mersions into submersions,  and also smooth objects into smooth ones
provided k' is projective as a module over k.    We also have a \uline{ring}
\uline{extension functor} from U(k) to U(k'),  giving rise to a diagram (I')
similar to (I) above,  it transforms submersions into submersions,
smooth objects into smooth ones.

The smoothness condition is likely to come in in two ways,  one
is via \uline{flatness}  (we may call an object of U(k) ``flat'' when it is iso-%
morphic to some W(M) with M a flat k-module),  whereas projectivity is
needed in order to insure that in certain cases,  weak equivalences are
homotopisms.  Flatness is the kind of condition which insures the vali-%
dity of ``naive'' universal coefficients formulae for homotopy or homology
objects ,  whereas projectivity may be needed in case of such formulae
for cohomology rather than homology.

\margin{31}
